% Generated draft: requires mathematical and source review.
\documentclass[11pt]{article}
\usepackage[margin=25mm]{geometry}
\usepackage{fontspec}
\setmainfont{Latin Modern Roman}
\usepackage{amsmath,amssymb,mathtools,unicode-math}
\setmathfont{Latin Modern Math}
\usepackage{enumitem,needspace,xurl,hyperref}
\setlength{\parindent}{0pt}
\setlength{\parskip}{5pt}
\setlength{\emergencystretch}{3em}
\newcommand{\N}{\mathbb{N}}
\newcommand{\Z}{\mathbb{Z}}
\newcommand{\Q}{\mathbb{Q}}
\newcommand{\R}{\mathbb{R}}
\newcommand{\C}{\mathbb{C}}
\newcommand{\F}{\mathbb{F}}
\newcommand{\A}{\mathbb{A}}
\newcommand{\PP}{\mathbb{P}}
\newcommand{\E}{\mathbb{E}}
\newcommand{\eps}{\varepsilon}
\newcommand{\dd}{\,\mathrm{d}}
\begin{document}
\section*{UnsolvedMath ID: 6700039}
\textbf{Source problem code:} catalogue entry
\subsection{Definitions and mathematical statement}
For a smooth compact manifold \(X\), let \(\operatorname{Met}(X)\) denote its smooth Riemannian metrics. Positive scalar curvature means
\[
\exists g\in\operatorname{Met}(X)\ \forall x\in X,
\qquad \operatorname{Sc}_g(x)>0.
\]
The question asks whether such an \(X\) can nevertheless have no area-extremal or length-extremal metric in the source's sense.

The familiar pointwise length-domination relation between metrics is
\[
h\succeq_{\mathrm{len}}g\iff h(v,v)\ge g(v,v)\text{ for every tangent vector }v.
\]
Area domination instead compares induced two-vector areas:
\[
|u\wedge v|_h^2=h(u,u)h(v,v)-h(u,v)^2
\ge |u\wedge v|_g^2
\]
for every tangent pair \(u,v\). These relations explain the comparison directions, but do not define extremal without specifying the admissible maps/metrics and scalar-curvature conclusion. Extremality and an additional equality-rigidity conclusion are also different conditions.

Writing \(\mathsf{AE}_X(h)\) and \(\mathsf{LE}_X(h)\) for the missing source predicates, the strong no-either-type reading is
\[
\exists X\quad
\bigl(\exists g\in\operatorname{Met}(X):\operatorname{Sc}_g>0\bigr)
\ \land\
\bigl(\forall h\in\operatorname{Met}(X),\
\neg\mathsf{AE}_X(h)\land\neg\mathsf{LE}_X(h)\bigr).
\]
The source does not say whether it intends this joint absence or separate existence questions for each extremality notion, or whether the last metric class is itself restricted to positive scalar curvature. No dimension, boundary convention, scalar-bound comparison, or map-degree/normalization data for extremality are supplied. Thus no maximization functional or rigidity theorem is invented from the phrase extremal.

The supplied research summary calls the existence question open and says no resolution was located. Its informal description as maximal area/length rigidity is not a precise definition of either predicate.
\subsection{Short English statement}
Can a compact manifold admit positive scalar curvature yet admit neither an area-extremal nor a length-extremal metric? The two extremality definitions and the scope of the metric quantifier must be specified.
\subsection{Sources}
\begin{itemize}
\item Local UnsolvedMath registry, permanent ID \(6700039\), exact source code \texttt{catalogue entry}; printed scalar-curvature tag [?41], occurrence 41. Proposer Misha Gromov, 2017.
\item Gromov, \emph{101 Questions, Problems and Conjectures around Scalar Curvature} (2017), recorded PDF page 48; \url{https://www.ihes.fr/~gromov/wp-content/uploads/2018/08/101-problemsOct1-2017.pdf}. Recorded access 2026-07-29; pdf-text-normalized extraction.
\item Supplied OPEN-TRIAGE research summary reports no located resolution, but supplies no precise extremality definition or additional publication citation.
\end{itemize}
\end{document}
